%% ============================================================
%% CodeChange20260519.tex — ARES 코드 수정 보고서
%% BLE 명령 파이프라인 최적화 (Web/commandexecutor.js, Pico/main.py)
%% XeLaTeX 컴파일: xelatex -shell-escape CodeChange20260519.tex
%% ============================================================
\documentclass[11pt, oneside]{article}

%% ----- 한국어 / 폰트 -----
\usepackage{kotex}
\usepackage{fontspec}
\defaultfontfeatures{Mapping=}
%% 사용 환경에 맞게 수정. macOS: /Users/<USER>/Library/Fonts/
\newcommand{\userfontpath}{C:/Users/young/AppData/Local/Microsoft/Windows/Fonts/}

\setmainfont{KoPub Batang}[
  Path=\userfontpath, Scale=0.92, Mapping=, LetterSpace=-2.0,
  UprightFont = KoPub Batang Medium.ttf,
  BoldFont    = KoPub Batang Bold.ttf
]
\setsansfont{KoPub Dotum}[
  Path=\userfontpath, Scale=0.92, Mapping=, LetterSpace=-2.0,
  UprightFont = KoPub Dotum Medium.ttf,
  BoldFont    = KoPub Dotum Bold.ttf
]
\setmonofont{D2Coding}[
  Path=\userfontpath, Scale=0.88,
  Extension = .ttc,
  UprightFont = D2Coding-Ver1.3.2-20180524-all
]
\setmainhangulfont{KoPub Batang}[
  Path=\userfontpath, Scale=0.92, Mapping=, LetterSpace=-2.0,
  UprightFont = KoPub Batang Medium.ttf,
  BoldFont    = KoPub Batang Bold.ttf
]
\setsanshangulfont{KoPub Dotum}[
  Path=\userfontpath, Scale=0.92, Mapping=, LetterSpace=-2.0,
  UprightFont = KoPub Dotum Medium.ttf,
  BoldFont    = KoPub Dotum Bold.ttf
]
\setmonohangulfont{KoPub Batang}[
  Path=\userfontpath, Scale=0.92, LetterSpace=-2.0,
  UprightFont = KoPub Batang Medium.ttf,
  BoldFont    = KoPub Batang Bold.ttf
]

%% ----- 레이아웃 -----
\usepackage{geometry}
\geometry{
  paperwidth=190mm, paperheight=240mm,
  top=24mm, bottom=28mm, inner=24mm, outer=22mm,
  headheight=14pt, headsep=8mm, footskip=28pt
}
\linespread{1.12}
\setlength{\parskip}{0.85em}
\setlength{\parindent}{0pt}
\raggedbottom
\setcounter{tocdepth}{2}
\setcounter{secnumdepth}{2}

%% ----- 색상 / 박스 / 코드 -----
\usepackage{xcolor}
\definecolor{deepblue}{RGB}{22, 56, 100}
\definecolor{accentcopper}{RGB}{200, 105, 50}
\definecolor{labcolor}{RGB}{40, 140, 125}
\definecolor{cautioncolor}{RGB}{195, 85, 70}
\definecolor{rulegray}{RGB}{210, 214, 220}
\definecolor{codebackground}{RGB}{249, 250, 252}

\usepackage{minted}
\usemintedstyle{friendly}
\setminted{
  fontsize=\footnotesize,
  bgcolor=codebackground,
  frame=leftline,
  framerule=0.8pt,
  rulecolor=\color{accentcopper},
  breaklines=true,
  breakanywhere=true,
  tabsize=2,
  linenos=false,
  autogobble=true,
}

\usepackage{tcolorbox}
\tcbuselibrary{skins, breakable}

\newtcolorbox{keypoint}{
  enhanced, breakable, sharp corners, arc=0pt, frame hidden, boxrule=0pt,
  colback=accentcopper!6,
  borderline west={2pt}{0pt}{accentcopper},
  top=6pt, bottom=6pt, left=12pt, right=10pt,
  before skip=1em, after skip=0.9em,
}
\newtcolorbox{caution}{
  enhanced, breakable, sharp corners, arc=0pt, frame hidden, boxrule=0pt,
  colback=cautioncolor!6,
  borderline west={2pt}{0pt}{cautioncolor},
  top=6pt, bottom=6pt, left=12pt, right=10pt,
  before skip=1em, after skip=0.9em,
}

%% ----- 표 -----
\usepackage{booktabs}
\usepackage{tabularx}
\usepackage{array}

%% ----- 섹션 스타일 (간결) -----
\usepackage{titlesec}
\titleformat{\section}[block]
  {\sffamily\bfseries\Large\color{deepblue}}{\thesection}{0.6em}{}
  [\vspace{2pt}{\color{rulegray}\hrule height 0.5pt}\vspace{4pt}]
\titleformat{\subsection}
  {\sffamily\bfseries\normalsize\color{accentcopper}}{\thesubsection}{0.6em}{}

%% ----- 헤더 / 푸터 -----
\usepackage{fancyhdr}
\pagestyle{fancy}
\fancyhf{}
\renewcommand{\headrulewidth}{0pt}
\renewcommand{\footrulewidth}{0pt}
\fancyfoot[L]{\sffamily\small\color{deepblue}\bfseries ARES BLE 파이프라인 최적화}
\fancyfoot[R]{\sffamily\small\thepage}

\usepackage{hyperref}
\hypersetup{
  colorlinks=true, linkcolor=deepblue, urlcolor=deepblue,
  pdftitle={ARES 코드 수정 보고서 — 2026-05-19},
  pdfauthor={(주)코리아사이언스, 동명대학교 게임학부},
}

\renewcommand{\contentsname}{목차}
\renewcommand{\figurename}{그림}
\renewcommand{\tablename}{표}

\title{%
  {\sffamily\bfseries\fontsize{20}{24}\selectfont\color{deepblue}ARES 코드 수정 보고서}\\[4pt]
  {\rmfamily\normalsize\color{accentcopper}BLE 명령 파이프라인 최적화}
}
\author{\sffamily (주)코리아사이언스\quad/\quad 동명대학교 게임학부}
\date{2026년 5월 19일}

\begin{document}
\maketitle
\thispagestyle{plain}

\section{개요}

ARES 로버의 블록 코딩 사용 중 ``LED 1번 켜기 → LED 2번 켜기''와 같이
연속된 명령 사이에 체감 가능한 지연이 발생하던 문제를 분석하여, 명령당
누적 약 250ms에 달하던 지연을 출력 명령 기준 약 40ms 수준으로 줄였다.
변경 대상은 \texttt{Web/commandexecutor.js}와 \texttt{Pico/main.py}이며,
명령의 의미에 따라 응답 라운드트립을 \emph{선택적으로} 생략하는 정책을
양쪽에 일관되게 도입하였다.

\section{문제 진단}

단일 LED 명령 한 번이 종단 간에 소요하는 시간 예산은 다음과 같이
구성되어 있었다.

\begin{table}[h]
\centering
\small
\begin{tabularx}{\linewidth}{@{}llXr@{}}
\toprule
\# & 위치 & 원인 & 시간 \\
\midrule
1 & \texttt{commandexecutor.js:199} & \texttt{sendData(cmd, true)}로 모든 명령에 응답 대기 & 라운드트립 \\
2 & \texttt{bluetooth.js:396} & \texttt{CHUNK\_DELAY=50ms} (\texttt{constants.js:14}) & 50ms \\
3 & \texttt{main.py:86} & \texttt{RECEIVE\_TIMEOUT\_MS=50ms} 무조건 대기 & 50ms \\
4 & \texttt{commandexecutor.js:213} & 모든 명령 후 \texttt{setTimeout(100)} & 100ms \\
5 & BLE GATT  & connection interval (HM-10 기본) & 30--70ms \\
\bottomrule
\end{tabularx}
\caption{단일 LED 명령의 누적 지연(개선 전) — 합계 약 250ms.}
\end{table}

\begin{keypoint}
정상 동작이라도 누적 지연이 \textbf{약 250ms / 명령}에 달했다. 카운트
다운, 파티 조명처럼 짧은 간격의 LED 시퀀스가 ``끊기는'' 체감의 주된
원인이다.
\end{keypoint}

\section{변경 정책 — 명령별 분리}

명령을 두 부류로 나누어 응답 대기 여부와 cooldown을 다르게 적용한다.

\begin{table}[h]
\centering
\small
\begin{tabularx}{\linewidth}{@{}lXl@{}}
\toprule
부류 & 명령 & 응답 대기 \\
\midrule
Fire-and-forget & \texttt{LED\_ON}, \texttt{LED\_OFF}, \texttt{MAIN\_LED\_*}, \texttt{[v0 v1 v2 v3 v4]} (LED 패턴),
\texttt{MSG}, \texttt{CLEAR\_DISPLAY}, \texttt{SERVO\_FORWARD}/\texttt{BACKWARD}/\texttt{LEFT}/\texttt{RIGHT}/\texttt{STOP} (연속),
\texttt{DC\_FORWARD}/\texttt{BACKWARD}/\texttt{STOP} (연속), \texttt{GUN\_FIRE} & 아니오 \\
\midrule
응답 대기 유지 & \texttt{BUZZER\_ON,F,D} (D초 blocking), \texttt{SERVO\_t*}/\texttt{DC\_t*} (N초 blocking),
\texttt{SLEEP,N} (N초 blocking), \texttt{DISTANCE}, \texttt{MAGNET}, \texttt{PING} & 예 \\
\bottomrule
\end{tabularx}
\caption{명령 부류와 응답 대기 정책. Pico에서 blocking으로 처리되거나
값을 반환하는 명령은 응답 대기를 유지하여 타이밍과 값 정합성을 보존한다.}
\end{table}

\section{Web 측 변경 — \texttt{Web/commandexecutor.js}}

\subsection{핵심 코드}

\begin{minted}{javascript}
// 즉시 완료 명령 — Pico가 blocking 처리하지 않으므로 응답 대기 불필요.
// 시간 의존적(BUZZER_ON, SERVO_t*, DC_t*, SLEEP) 또는 값 반환(DISTANCE,
// MAGNET, PING)은 이 집합에 넣지 말 것. 응답 대기로 두어야 타이밍/값이
// 보장된다.
FIRE_AND_FORGET_HEADS: new Set([
  'LED_ON', 'LED_OFF',
  'MAIN_LED_ON', 'MAIN_LED_OFF',
  'MSG', 'CLEAR_DISPLAY',
  'SERVO_FORWARD', 'SERVO_BACKWARD', 'SERVO_LEFT', 'SERVO_RIGHT', 'SERVO_STOP',
  'DC_FORWARD', 'DC_BACKWARD', 'DC_STOP',
  'GUN_FIRE',
]),

_isFireAndForget(command) {
  if (command.startsWith('[')) return true;     // LED 패턴 [v0 v1 v2 v3 v4]
  const head = command.split(',')[0];
  return this.FIRE_AND_FORGET_HEADS.has(head);
},

async sendCommand(command) {
  if (!state.isExecuting) {
    Logger.add('[중단] 실행이 중단되었습니다', 'warning');
    return;
  }
  BluetoothManager.updateStatus('명령 실행 중...', STATUS_COLORS.ORANGE);

  const fireAndForget = this._isFireAndForget(command);

  try {
    await BluetoothManager.sendData(command, !fireAndForget);
    if (DEBUG) Logger.add(`[완료] ${command}`, 'info');
  } catch (error) {
    // ... 에러 처리 ...
  }

  // 송신 간격 가드.
  // - fire-and-forget: UART(9600 baud)/BLE 버퍼 오버플로우 방지 최소 마진.
  // - 응답 대기: 이미 라운드트립을 거쳤으므로 가드 단축.
  const cooldown = fireAndForget ? 40 : 20;
  await new Promise(resolve => setTimeout(resolve, cooldown));
}
\end{minted}

\subsection{효과}

\begin{itemize}
  \item 출력 명령(LED, SERVO 연속, DC 연속 등): BLE 라운드트립과 100ms
        cooldown이 모두 사라짐 → 명령당 약 40ms.
  \item 응답 대기 명령: cooldown만 100ms→20ms 단축. 의미적 동기화는
        그대로 유지.
\end{itemize}

\section{Pico 측 변경 — \texttt{Pico/main.py}}

웹의 분류와 \emph{동일 기준}으로 Pico에서도 응답 송신 여부를
결정한다. 또한 단일 chunk 명령에 대한 무조건 50ms 대기를 제거한다.

\subsection{NO\_RESPONSE\_CMDS 및 \texttt{\_needs\_response}}

\begin{minted}{python}
class AresRover:
    # 응답 송신 생략 명령 (Web/commandexecutor.js의 FIRE_AND_FORGET_HEADS와 동기화).
    # 즉시 완료되는 출력 명령은 응답 라운드트립을 생략하여 처리량을 높이고,
    # 응답 매칭 어긋남(LED_ON 응답이 다음 명령 슬롯에 잘못 들어가는 현상)을 차단한다.
    NO_RESPONSE_CMDS = (
        "LED_ON", "LED_OFF", "MAIN_LED_ON", "MAIN_LED_OFF",
        "MSG", "CLEAR_DISPLAY",
        "SERVO_FORWARD", "SERVO_BACKWARD", "SERVO_LEFT", "SERVO_RIGHT", "SERVO_STOP",
        "DC_FORWARD", "DC_BACKWARD", "DC_STOP",
        "GUN_FIRE",
    )

    def _needs_response(self, data):
        """응답 송신 여부. fire-and-forget 명령은 False (NO_RESPONSE_CMDS 참조)."""
        if data.startswith("["):           # LED 패턴 [v0 v1 v2 v3 v4]
            return False
        head = data.split(",", 1)[0]
        return head not in self.NO_RESPONSE_CMDS
\end{minted}

\subsection{\texttt{\_read\_uart\_line()} 폴링 방식 교체}

기존 구현은 매 호출마다 \texttt{utime.sleep\_ms(RECEIVE\_TIMEOUT\_MS)}로
50ms를 무조건 대기한 뒤 두 번째 read를 수행했다. newline이 이미 도착한
단일 chunk 명령에도 적용되어 처리량 천장이 약 16--17 cmd/s에 묶여
있었다. 폴링으로 교체하여 newline 발견 즉시 종료하도록 한다.

\begin{minted}{python}
def _read_uart_line(self):
    """UART에서 완전한 라인 읽기.
    newline 발견 즉시 종료 — 단일 chunk 명령에서 무조건 50ms를 기다리지 않는다.
    멀티 chunk 명령(LED 패턴, SYS_SET 등)은 RECEIVE_TIMEOUT_MS까지 폴링하여 안전.
    """
    if not self.uart.any() and not self.rx_buffer:
        return None

    start = utime.ticks_ms()
    while utime.ticks_diff(utime.ticks_ms(), start) < RECEIVE_TIMEOUT_MS:
        if self.uart.any():
            chunk = self.uart.read()
            if chunk:
                self.rx_buffer += chunk.decode('utf-8', 'ignore')

            # 버퍼 오버플로우 방지
            if len(self.rx_buffer) > MAX_BUFFER_SIZE:
                print("[UART] 버퍼 오버플로우, 초기화")
                self.rx_buffer = ""
                return None

            # newline 즉시 종료
            if '\n' in self.rx_buffer:
                newline_idx = self.rx_buffer.index('\n')
                complete_line = self.rx_buffer[:newline_idx].strip()
                self.rx_buffer = self.rx_buffer[newline_idx + 1:]
                return complete_line
        else:
            utime.sleep_ms(2)

    # 타임아웃: 라인 종료자 없는 부분 데이터 반환
    if self.rx_buffer:
        data = self.rx_buffer.strip()
        if data and not data.endswith(','):
            self.rx_buffer = ""
            return data

    return None
\end{minted}

\subsection{\texttt{run()}의 응답 조건부 송신}

\begin{minted}{python}
# 변경 전
else:
    response = self._process_command(data)
    self._send_response(response)

# 변경 후
else:
    response = self._process_command(data)
    if self._needs_response(data):
        self._send_response(response)
\end{minted}

\section{기대 효과 (정량)}

\begin{table}[h]
\centering
\small
\begin{tabular}{@{}lcc@{}}
\toprule
명령 종류 & 변경 전 & 변경 후 \\
\midrule
\texttt{LED\_ON} 등 fire-and-forget & 약 250ms & 약 40ms \\
\texttt{BUZZER\_ON,F,D} 등 blocking & $D + 250$ms & $D + 20$ms \\
\texttt{DISTANCE}/\texttt{MAGNET} 등 값 반환 & 약 250ms & RTT $+$ 20ms \\
Pico 처리량(단일 chunk) & $\approx 16$--$17$ cmd/s & $\approx 100^+$ cmd/s \\
\bottomrule
\end{tabular}
\caption{명령 부류별 종단 간 지연 비교 (단위: ms).}
\end{table}

추가로 응답 매칭 어긋남(이전 fire-and-forget 명령의 응답이 다음
응답 대기 명령의 \texttt{pendingResolve} 슬롯에 잘못 들어가는 코너
케이스) 위험이 \emph{구조적으로} 제거된다. Pico에서 응답 자체를
보내지 않으므로 침범할 응답이 존재하지 않는다.

\section{잔여 위험 및 후속 작업}

\begin{caution}
변경된 두 파일은 \textbf{같이 갱신}되어야 한다. 한쪽만 갱신되면:
\begin{itemize}
  \item Web만 갱신: 출력 명령에 응답 대기를 안 하지만 Pico는 여전히
        응답을 보냄 → BLE notify 누적, 응답 매칭 어긋남 위험 잔존.
  \item Pico만 갱신: Web이 출력 명령에 응답 대기를 하지만 Pico가 응답을
        안 보냄 → 매 명령 5초 timeout.
\end{itemize}
\end{caution}

\subsection{남은 병목 — UART baud rate}

가장 본질적인 천장은 BLE 모듈 ↔ Pico 사이의 9600 baud UART이다.
HM-10/BT05의 AT 명령(\texttt{AT+BAUD8} 등)으로 38400 또는 115200 baud로
상향하고 \texttt{Pico/main.py}의 \texttt{UART\_BAUDRATE}를 동기화하면
byte당 전송 시간을 4--12배 단축할 수 있다. BLE 모듈 펌웨어 설정과
Pico 코드 동시 수정이 필요하므로 별도 작업으로 분리한다.

\subsection{권장 회귀 테스트}

\begin{enumerate}
  \item \textbf{카운트다운} — \texttt{LED\_OFF} 5개 연속 소등. 명령
        사이 체감 지연이 거의 사라져야 한다.
  \item \textbf{LED 패턴 + 거리 센서 혼합} — \texttt{[v0 v1 v2 v3 v4]}
        직후 \texttt{DISTANCE}. 거리값이 \texttt{1} 같은 LED 응답으로
        오염되지 않는지 확인.
  \item \textbf{부저 사이렌} — \texttt{BUZZER\_ON,400,0.5} 와
        \texttt{BUZZER\_ON,800,0.5}의 반복. 타이밍 정확성이 유지되어야
        한다.
  \item \textbf{시간 제한 모터} — \texttt{SERVO\_tFORWARD,N},
        \texttt{DC\_tFORWARD,N}의 N초 동기화가 그대로 유지되는지 확인.
\end{enumerate}

\section{변경 파일 요약}

\begin{table}[h]
\centering
\small
\begin{tabularx}{\linewidth}{@{}lX@{}}
\toprule
파일 & 변경 요지 \\
\midrule
\texttt{Web/commandexecutor.js} & \texttt{FIRE\_AND\_FORGET\_HEADS} 집합 도입, \texttt{sendCommand()}에서 분기 처리, cooldown 분리(40 / 20ms). \\
\texttt{Pico/main.py} & \texttt{NO\_RESPONSE\_CMDS} 클래스 변수, \texttt{\_read\_uart\_line()} 폴링 방식으로 교체, \texttt{\_needs\_response()} 도입, \texttt{run()}에서 조건부 \texttt{\_send\_response()}. \\
\bottomrule
\end{tabularx}
\caption{본 보고서가 다루는 코드 수정 범위.}
\end{table}

\end{document}
